\chapter{Introduction}\label{chap:intro}

\section{Motivation}

Unlike screen-based assistants, social robots occupy the same physical space as their users. This embodiment changes what users expect from them. A person addressing a robot in a room may reasonably expect it to know what is nearby, who is present, and which objects are being discussed.

Pepper is a humanoid social robot built around speech, expressive motion, camera input, tablet output, and perception services provided through the NAOqi platform. These capabilities make it suitable for situated dialogue, but they also expose a central limitation: a robot that speaks about its environment should be able to connect its utterances to what it perceives. In this thesis, Pepper is treated as a situated communication partner that answers questions about its surroundings.

Recent Pepper systems show that this direction is technically and practically relevant. Pepper has been used as a teaching assistant and as a service robot, for example in library assistance with visual perception and OCR~\cite{sievers_humanoid_2025,trinquet_future_2025}. Other systems connect Pepper or similar robots to external language, vision, or retrieval components to support grounded dialogue~\cite{brad_retrieval-augmented_2025,grassi_grounding_2024,latif_physicsassistant_2024}. These works show how modern AI components can extend situated interaction. The central problem of this thesis is therefore not simply to attach a language model to Pepper, but to ground its answers in an explicit, dynamically updated model of the perceived scene.

Large language models (LLMs) make this goal more plausible. They can make robotic dialogue more flexible and natural than hand-written response templates or rigid dialogue states defined by dialogue management systems~\cite{aturhurra_leveraging_2024}. However, an LLM does not automatically know what Pepper currently sees. If the robot is asked about a laptop or a chair placed in front of it, the answer must be supported by perception and memory rather than by the model's general language knowledge.

Object detection alone is too thin a basis for this kind of dialogue. A detector can localize objects and return labels and bounding boxes, but many questions require relations: what is beside the laptop, what is on the table, or who is looking toward the robot. Scene graphs provide a structured layer between low-level perception and high-level dialogue by encoding objects, their attributes, and their relations~\cite{johnson_image_2015}, which makes them suitable for conversational grounding.

Scene awareness also requires persistence. Real interaction is not limited to one image: Pepper may inspect a room, answer a question, move its head, observe another part of the scene, and later be asked about an object from an earlier view. The system therefore needs scene memory: it must preserve object identities across observations, update attributes and relations, and expose the resulting state to the dialogue layer.

Pepper also imposes a technical constraint. Modern perception and language models are too demanding to run as a full pipeline on Pepper's onboard hardware. In this thesis, Pepper is therefore used as an embodied client, while an external machine performs the computationally expensive perception, memory, and language processing. This follows the same practical direction as previous Pepper systems that rely on external computation for modern perception or dialogue capabilities~\cite{reyes_near_2018,trinquet_future_2025,brad_retrieval-augmented_2025,grassi_grounding_2024}.

The target setting is an ordinary indoor environment, such as an office, classroom, or computer laboratory. In this setting, the robot should be able to observe its surroundings, build a structured memory of relevant objects and relations, and use this memory when answering the user.

\section{Problem Statement and Objectives}

The central problem addressed in this thesis is how to enable the Pepper robot to conduct dialogue that is grounded in perception and persistent across observations. A scene-aware answer should not depend only on general language-model knowledge; it should also reflect what the robot has observed, which objects and people are present, and how these entities are related.

The thesis is system-oriented. It does not propose a new object detector, scene graph model, or language-model architecture. Instead, it investigates how existing perception, vision-language, and language components can be selected, combined, and adapted for a Pepper-based interaction pipeline.

The main objective is to design and implement a distributed architecture in which Pepper serves as the embodied interaction client, while an external server performs the demanding inference and state management. Within this architecture, the system detects relevant objects and people, preserves object identities across discontinuous observations, organizes perceived entities into structured scene memory, and uses that memory as grounding context for dialogue.

\section{Main Contributions}

The contribution of this thesis is the design, implementation, and evaluation of a complete scene-aware dialogue system for the Pepper robot. The system integrates robot-side interaction, server-side multimodal perception, persistent memory, and grounded language generation into a single pipeline.

At the representation level, the thesis introduces a perception-to-memory pipeline that detects objects and people, supports identity persistence across discontinuous observations, and stores the perceived environment in a form suitable for dialogue.


At the interaction level, it implements a dialogue layer that conditions answers on the stored scene state rather than on language-model knowledge alone.


Finally, the evaluation covers graph quality, Set-of-Mark prompting, vocabulary design, prompt formatting, and robot-side interaction latency.

\section{Thesis Structure}

Chapter~\ref{chap:background} reviews the background and related work needed for the thesis, focusing on HRI, Pepper, visual perception, scene graphs, vision-language models, large language models, and grounded dialogue. Chapter~\ref{chap:system_architecture_rewrite} develops the architecture and design of the proposed system. Chapters~\ref{chap:server_side_rewrite} and~\ref{chap:robot_side_rewrite} detail the server-side and robot-side implementation. Chapter~\ref{chap:experiment} reports the testing and evaluation, including scene graph generation, Set-of-Mark prompting, vocabulary sensitivity, prompt formatting, and latency. Chapter~\ref{chap:end} concludes with the results, limitations, and possible future work.
